\documentclass[10pt,landscape,letterpaper]{cheatsheet}

\usepackage{amsmath,amsthm,amssymb}
\usepackage{physics}

\newcommand{\PRLsep}{\noindent\makebox[\linewidth]{\resizebox{0.3333\linewidth}{1pt}{$\bullet$}}\bigskip}
\newcommand{\laplace}[1]{\mathcal{L}\left\{#1\right\}}
\newcommand{\laplacei}[1]{\mathcal{L}^{-1}\left\{#1\right\}}
\DeclareMathOperator{\erfc}{erfc}
\newcommand{\highlight}[1]{\colorbox{yellow}{$\displaystyle #1$}}

\title{Methods Cheat Sheet}
\author{Jeremy Favro}
\date{\today\\Revision 1}

\begin{document}

\maketitle

\section{Power Series}
For an ODE of the form 
$y''+p(x)y'+q(x)y=0$, the solution is of the form
$\sum_{n=0}^{\infty}a_nx(x-x_0)^n$ if $p,q$ are analytic at $x_0$. This power
series is guaranteed to converge for $\abs{x-x_0}<R$ where $R$ is the distance from $x_0$ to the nearest singular point.
\subsection{Common Power Series}
$e^x=\sum_{n=0}^{\infty}\frac{x^n}{n!}$\\
$\cos x = \sum_{n=0}^{\infty}(-1)^n\frac{x^{2n}}{(2n)!}$\\
$\sin x = \sum_{n=0}^{\infty}(-1)^n\frac{x^{2n+1}}{(2n+1)!}$\\
$\cosh x = \sum_{n=0}^{\infty}\frac{x^{2n}}{(2n)!}$\\
$\sinh x = \sum_{n=0}^{\infty}\frac{x^{2n+1}}{(2n+1)!}$\\
$\frac{1}{1-x} = \sum_{n=0}^{\infty}x^n$\\
$\frac{1}{1+x} = \sum_{n=0}^{\infty}(-x)^n$\\
$\tan^{-1} x = \sum_{n=0}^{\infty}(-1)^n\frac{x^{2n+1}}{2n+1}$\\
$\ln(1+x) x = \sum_{n=0}^{\infty}(-1)^{n+1}\frac{x^n}{n}$\\
\subsection{Ordinary Points}
A point is called ordinary if $p$ and $q$ are analytic at that point.
\subsection{Solutions Around Singular Points}
A point about which $p,q$ are \emph{not analytic} but 
$(x-x_0)p(x)$ and $(x-x_0)^2q(x)$ are is called regular singular. 
We are guaranteed one solution of the form 
$\sum_{n=0}^{\infty}a_n(x-x_0)^{n+r}$ where 
$r$ is determined using the \emph{indicial equation}
$r(r-1)+\left[\lim_{x \to x_0}(x-x_0)p(x)  \right]r+\left[\lim_{x \to x_0}(x-x_0)^2q(x)  \right]=0$. We always use the larger $r$, $r_1$ first.
This series is guaranteed to converge for $0<x-x_0<R$. If the difference between the roots, $r_1-r_2$ is not an integer
then the series gives both linearly independent solutions. If the difference is an integer then the series gives at least one solution, with the other either
give directly or given by $Cy_1\ln(x)+\sum_{n=0}^{\infty}b_n(x-x_0)^{n+r_2}$. When the roots are repeated our second solution is guaranteed to be 
$y_1\ln(x)+\sum_{n=0}^{\infty}b_n(x-x_0)^{n+r_2}$
\section{Special Functions}
\subsection{Elementary Functions}
\subsection{Gamma Function}
$\Gamma(p)=\int_{0}^{\infty}x^{p-1}e^{-x}\,dx$\\
$\Gamma(n+1)=n!\forall n\in\mathbb{N}$\\
$\Gamma(p+1)=p\Gamma(p)$\\
$\Gamma(1/2)=\sqrt{\pi}$\\
$\Gamma(p)\Gamma(1-p)=\frac{\pi}{\sin(\pi p)}0<p<1$\\
$\Gamma(p)\Gamma(1-p)=\int_{0}^{\infty}\frac{x^{p-1}}{1+x}\,dx\,0<p<1$\\
$\Gamma(p)\approx\sqrt{\frac{2\pi}{p}}\left(\frac{p}{e}\right)^p,p\to\infty$
\subsection{Beta Function}
$\beta(p,q)=\beta(q,p)=\int_{0}^{1}x^{p-1}(1-x)^{q-1}\,dx\,p,q>0$\\
$\beta(p,q)=\frac{1}{({b-a})^{p+q-1}}\int_{a}^{b}(x-a)^{p-1}(b-x)^{q-1}$\\
$\beta(p,q)=2\int_{0}^{\pi/2}\sin^{2p-1}(\theta)\cos^{2q-1}(\theta)d\theta$\\
$\beta(p,q)=\frac{\Gamma(p)\Gamma(q)}{\Gamma(p+q)}$
\subsection{Error Function}
$\erf(x)=\frac{2}{\sqrt{\pi}}\int_{0}^{x}e^{-t^2}dt$\\
$\erfc(x)=1-\erf(x)$\\
$\erf(0)=0$\\
For $x\ll 1$, $\erf(x)\approx \frac{2}{\sqrt{\pi}}\left[x-\frac{x^3}{3}+\frac{x^5}{2!5}-\frac{x^7}{3!7}+\dots\right]$\\
For large $x$,
$\erf(x)\approx\frac{e^{-x^2}}{x\sqrt{\pi}}\left[1-\frac{1}{2x^2}+\frac{1\cdot 3}{(2x^2)^2}+\dots\right]$.
\subsection{Legendre Polynomials}
ODEs of the form 
$(1-x^2)y''-2xy'+\ell(\ell+1)y=0$
are solved by
$c_1Q_\ell(x)+c_2P_\ell(x)$.\\
$P_\ell(x)=\frac{1}{2^\ell}\sum_{m=0}^M\frac{(-1)^m(2\ell-2m)!}{m!(\ell-m)!(\ell-2m)!}$\\
where
$M=\begin{cases}
    \ell/2&\ell\text{ even}\\
    (\ell-1)/2&\ell\text{ odd}
\end{cases}$
$P_\ell(x)=\frac{1}{2^\ell\ell!}\frac{d^\ell}{dx^\ell}\left(x^2-1\right)^\ell$
\subsubsection{Generating Function}
Useful in proving recurrence relations
$\Phi(x,h)=\left(1-2xh+h^2\right)^{-1/2}=\sum_{\ell=0}^{\infty}h^\ell P_\ell(x)$
\subsubsection{Recurrence Relations}
$(\ell+1)P_{\ell+1}(x)=(2\ell+1)xP_\ell(x)-\ell P_{\ell-1}(x)$\\
$xP_\ell'(x)-P_{\ell-1}'(x)=\ell P_\ell(x)$\\
$P_{\ell+1}'(x) -xP_\ell'(x)=(\ell+1)P_\ell(x)$\\
$(2\ell+1)P_\ell(x)=P_{\ell+1}'(x)-P_{\ell-1}'(x)$\\
$(1-x^2)P_\ell'(x)=\ell P_{\ell-1}(x)-\ell x P_\ell(x)$
\subsubsection{Orthogonality}
$\int_{-1}^{1}P_n(x)P_m(x)=\frac{2}{2n+1}\delta_{nm}$\\

\subsection{Bessel Functions}
ODEs of the form \\
$x^2y''+xy'+(\lambda^2x^2-\nu^2)y=0$
are solved by
$y=c_1J_\nu(\lambda x)+c_2Y_\nu(\lambda x)$.\\
$J_\nu(x)=\sum_{m=0}^\infty\frac{(-1)^mx^{2m+\nu}}{2^{2m+\nu}\Gamma(m+\nu+1)}$\\
$Y_\nu(x)=J_\nu(x)\ln(x) + \sum_{m=0}^\infty \frac{(-1)^m}{2^{2m}(m!)^2}x^{m+\nu}\cdot\sum_{n=1}^m\frac{1}{n}$.\\
ODEs of the form \\
$x^2y''+xy'-(\lambda^2x^2+\nu^2)y=0$
are solved by
$y=c_1I_\nu(\lambda x)+c_2K_\nu(\lambda x)$.\\
$I_\nu(x)=i^{-\nu}J_\nu(ix)$\\
$K_\nu(x)=\frac{\pi}{2}\frac{I_{-\nu}(x)-I_\nu(x)}{\sin(\nu\pi)}$. $I$ diverges for large $x$ and $K$ diverges for $x\to 0$.
ODEs of the form \\
$x^2y''+(1-2a)xy'+(b^2c^2x^{2c}+(a^2-p^2c^2))y=0$
are solved by
$y=c_1x^aJ_p(bx^c)+c_2x^aY_p(bx^c)$.\\
\subsubsection{Generating Function}
Useful in proving recurrence relations
$\Psi(x,h)=e^{\frac{1}{2}x(h-1/h)}=\sum_{n=-\infty}^{\infty}h^n J_n(x)$
\subsubsection{Properties}
$\frac{d}{dx}(x^\nu J_\nu(x))=x^\nu J_{\nu-1}(x)$\\
$\frac{d}{dx}(x^{-\nu} J_\nu(x))=-x^{-\nu} J_{\nu+1}(x)$\\
$J_{\nu-1}(x)+J_{\nu+1}=\frac{2\nu}{x}J_\nu$\\
$J_{\nu}'=\frac{1}{2}\left[J_{\nu-1}(x)+J_{\nu+1}\right]$\\
$xJ_{\nu}'=\nu J_{\nu}(x)-xJ_{\nu+1}$\\
$J_{1/2}(x)=\sqrt{\frac{2}{\pi x}}\sin(x)$\\
$J_{-1/2}(x)=\sqrt{\frac{2}{\pi x}}\cos(x)$\\
$J_{3/2}(x)=\sqrt{\frac{2}{\pi x}}\left[\frac{\sin(x)}{x}-\cos(x)\right]$\\
$J_{-3/2}(x)=-\sqrt{\frac{2}{\pi x}}\left[\frac{\cos(x)}{x}+\sin(x)\right]$
\subsubsection{Orthogonality}
$\int_0^1xJ_\nu(a_nx)J_\nu(b_mx)dx=\frac{1}{2}(J_\nu'(a))^2\delta_{nm}$
where $a_n$ and $b_m$ are the $n$th and $m$th zeroes of $J_\nu$.
\section{Self-Adjoint Operators and Sturm-Liouville}
A Sturm-Liouville problem is given by 
$L\left[\phi\right]+\lambda w(x)\phi(x)=0$. 
$L$ is a self-adjoint linear differential operator. For 
$L[\phi]=P_0\phi_{xx}+P_1\phi_x+P_2\phi$ this means that 
$P_0^\prime=P_1$. Distinct $\phi_n$ are orthogonal with respect to $w(x)$.
The coefficients of the eigenfunction expansion for $f(x)$ are 
$a_n=\int_a^bw(x)f(x)\phi_n(x)\,dx/\int_a^bw(x)\phi_n^2(x)\,dx$. We can transform a 
non-self-adjoint $L$ to a self adjoint one by multiplying by 
$P_0^{-1}\exp(\int P_1(x)/P_0(x)\,dx)$
\section*{Fourier Series}
For some function $f(x)$ periodic with period $2L$:\\
$a_0=\frac{1}{2L}\int_{-L}^{L}f(x)\,dx$\\
$a_n=\frac{1}{L}\int_{-L}^{L}f(x)\cos\left(\dfrac{n\pi x}{L}\right)\,dx$\\
$b_n=\frac{1}{L}\int_{-L}^{L}f(x)\sin\left(\dfrac{n\pi x}{L}\right)\,dx$\\
and
$f(x)=a_0+\sum_{n=1}^{\infty}a_n\cos\frac{n\pi x}{L}+b_n\sin\frac{n\pi x}{L}$
\end{document}